\chapter{結果分析與設計應用討論}
\label{chapter:results}

\section{整合系統設計流程回顧}
\label{sec:system_review}

第三、四章分別建立了 PI-ST-GNN 之方法論架構與訓練驗證結果。本章聚焦於整合系統在\textbf{實際設計情境}下的應用操作成果，包含：（1）多目標最佳化帕累托前緣之設計方案解讀，涵蓋熱舒適人行動線設計目標；（2）第~\ref{sec:building_extraction}節真實建物擷取流程於既有場地案例之驗證表現；以及（3）系統整合架構之效能瓶頸與最佳化方向討論。


\section{多目標最佳化設計方案分析}
\label{sec:pareto_analysis}

\subsection{NSGA-II 收斂過程與帕累托前緣形成}
\label{subsec:nsga2_convergence}

以訓練完成之 PI-ST-GNN 作為評估代理，NSGA-II 在 $80 \times 80$ 公尺場地尺度下同時最佳化第~\ref{subsec:design_variables}節定義之三項目標：$f_1$（最小化全場域時空平均 UTCI）、$f_2$（最小化人行動線時空平均熱暴露 $\bar{\Phi}_{\text{walkway}}$）與 $f_3$（最大化綠化率）。演算法設定如下：

\begin{itemize}
    \item \textbf{族群大小：}50 個個體
    \item \textbf{演化代數：}100 代
    \item \textbf{交叉算子：}模擬二進位交叉（SBX，交叉概率 0.9，$\eta_c = 15$）
    \item \textbf{突變算子：}多項式突變（突變概率 $1/n_{\text{var}}$，$\eta_m = 20$）
    \item \textbf{選擇策略：}可行性優先 + 非支配排序 + 擁擠距離選擇
\end{itemize}

由於每次個體評估之 GNN 推理已能提供立即性回饋（詳細延遲數值見附錄~\ref{appx:latency_breakdown}），100 代 $\times$ 50 個體規模的完整最佳化過程理論上可在數分鐘內完成，相較第~\ref{subsec:wallacei_comparison}節所述、直接呼叫高擬真度模擬引擎之 Wallacei 案例常需數週至數月的收斂時間，此一運算時間尺度的壓縮，是本研究得以在方案概念階段即進行多目標搜尋、而非僅能於方案定案後事後驗證的關鍵前提。

在正式導入第三項目標 $f_2$ 之前，本研究首先以一次僅含 $f_1$、$f_3$ 二目標之基準實驗（族群 40、演化 50 代、亂數種子 42，V5-300 checkpoint，驗證 $R^2=0.9837$），逐代真實呼叫 GNN 推理，於 GPU（NVIDIA RTX 5070 Ti）上共 404.2 秒完成，結果如圖~\ref{fig:nsga2_real}(a) 所示：族群最佳全域平均 UTCI 由第 1 代之 $36.29^\circ\text{C}$ 逐步下降至第 50 代之 $36.25^\circ\text{C}$，族群最佳綠化率則由 $0.009$ 上升至 $0.048$，可行解比例由 38/40 提升並穩定於 40/40。此一基準實驗之帕累托前緣形態，將於第~\ref{subsec:pareto_morphology}節詳述——與第~\ref{sec:v4_validation}節所述動態邊缺陷尚未修正之早期版本模型（V1--V4）之基準實驗結果不同，本次以 V5-300 執行之結果並未出現退化解現象，其原因與意義一併於下節討論。

\begin{figure}[htbp]
    \centering
    \includegraphics[width=\textwidth]{../urban-thermal-gnn/04_training/figures/fig_nsga2_convergence.png}
    \caption[基準二目標 NSGA-II 收斂結果]{基準二目標 NSGA-II 收斂結果（PI-ST-GNN V5-300 checkpoint 作為適應度求值器，族群 40、演化 50 代）}
    \label{fig:nsga2_real}
\end{figure}

上圖 (a) 為逐代族群最佳全域 UTCI（紅）與最佳綠化率（綠）之收斂曲線；(b) 為第 50 代最終帕累托前緣（$n=8$ 非支配解），色階代表各解之容積率（FAR）。


\subsection{帕累托前緣之真實形態與現行目標函數之侷限}
\label{subsec:pareto_morphology}

圖~\ref{fig:nsga2_real}(b) 之最終帕累托前緣揭示一項與本研究先前（V1--V4，動態邊缺陷尚未修正）基準實驗明顯不同之發現：6 個非支配解之 FAR 值收斂於 $2.03$--$2.05$ 之合理中高密度區間，落於第~\ref{subsec:coverage}節真實場景訓練集 FAR 分布之中段範圍內，\textbf{並未}出現先前版本模型所觀察到之「幾乎不興建任何建築量體」退化解（Degenerate Solution）現象。

此一對比之根本原因，最可能與第~\ref{sec:v4_validation}節、附錄~\ref{appx:dynamic_edges_fix}所揭露之動態邊缺陷有關：V1--V4 因動態邊快取從未真正建立，模型從未透過 \texttt{shadow}、\texttt{veg\_et} 兩類關係型邊實際學習建物遮蔭與植栽降溫之因果效應——即便 $f_1$（全域平均 UTCI）目標理論上應誘發模型辨識建物與喬木之降溫貢獻，一個從未見過真實遮蔭／降溫訊號傳遞路徑的模型，自然難以在其預測輸出中忠實反映此一效應，使最佳化引擎在缺乏正向誘因之情況下傾向「不蓋房子、盡量種樹」之退化解；V5-300 修正此一缺陷後，模型得以透過真實建立之動態邊學習建物與喬木對周邊空氣節點之實際降溫效果，NSGA-II 因而在同一二目標設定下即能辨識「適度建築量體可提供遮蔭、降低全域平均 UTCI」之正向訊號，收斂至合理密度區間而非退化解。換言之，先前版本所觀察到之退化解，其成因較可能主要並非二目標函數設計本身之結構性缺陷，而是動態邊缺陷導致模型未能學習建物／植栽之真實降溫因果關係；此一發現同時佐證了異質圖顯式邊結構（相對於純黑盒模型）於此類應用之關鍵價值——邊結構缺陷可被直接追溯並修正，且修正後之因果學習效果可於下游最佳化任務之行為變化中被觀察與驗證，此為本研究可解釋性設計（第~\ref{subsec:framework_interpretability}節）之具體效益展示。

儘管如此，現有 5 維違規向量（FAR、BCR、縮退、容納性、碰撞）仍僅檢核上限，未設定最低開發強度之下限約束，此一結構性缺口並未因動態邊修正而消失，僅是在 V5-300 修正後之模型行為下未被觸發；為避免此一風險於其他場地條件、氣候情境或搜索設定下重新浮現，第~\ref{sec:limitations}節與第~\ref{sec:future_work}節仍建議將最低容積率或建築棟數下限之可行性約束，列為後續強化措施。

此一發現對後續研究仍有直接啟示：第~\ref{subsec:from_astar_to_walkway}節與附錄~\ref{appx:walkway_derivation}所定義之第三項目標 $f_2$（人行動線熱暴露 $\bar{\Phi}_{\text{walkway}}$）之數學形式雖已完整推導，其研究動機——全空間平均指標難以反映行人沿路徑之動態熱暴露歷程——獨立於退化解問題而成立（第~\ref{sec:problem}節第二項研究缺口）；第~\ref{subsec:abc_comparison}節即以真實三目標實驗，驗證 $f_2$ 整合後之路徑感知評估能力與其對帕累托前緣量體分布之影響，並非作為退化解之修正手段，而是本研究對「路徑尺度熱舒適評估」此一研究缺口之直接回應。


\subsection{三目標實驗：路徑配置之三方案比較}
\label{subsec:abc_comparison}

為將第三項目標 $f_2$（人行動線熱暴露）實際整合進最佳化流程並驗證其路徑感知評估能力，本研究將 $f_2$ 正式接入 \texttt{FitnessEvaluator}（第~\ref{sec:nsga2}節），並針對同一 $80\times80$ 公尺場地與建築演化搜索空間，設計 A、B、C 三種固定人行路線配置，各自執行一次完整之三目標 NSGA-II 最佳化（族群 40、演化 50 代、亂數種子 42，V5-300 checkpoint，驗證 $R^2=0.9837$）：路線 A 沿建築臨街面配置、路線 B 貫穿開放廣場／中庭、路線 C 對角線橫越基地。此一設計直接回應口試委員對「建築形態與人行路徑之互動」須有具體比較實驗、而非僅呈現單一最佳化結果之要求。

三次獨立實驗之結果如圖~\ref{fig:abc_comparison} 所示，可歸納三項具體發現。其一，\textbf{三目標帕累托前緣呈現更寬廣之量體權衡區間}：三組實驗最終帕累托前緣之 FAR 範圍分別為路線 A（1.82--2.14）、路線 B（0.62--1.78）、路線 C（1.75--2.01），與第~\ref{subsec:pareto_morphology}節二目標基準實驗收斂於單一 FAR 區間（2.03--2.05）相比，三目標版本之非支配解涵蓋明顯更寬之量體區間，尤其路線 B 之帕累托前緣向下延伸至 FAR 0.62，顯示人行動線熱暴露目標確實為最佳化引擎引入了額外的量體—路徑權衡維度：部分非支配解以降低建築量體換取路徑熱舒適之改善，此一權衡在僅有 $f_1$、$f_3$ 之二目標設定下並不存在。三組實驗之 FAR 範圍均落於第~\ref{subsec:coverage}節真實場景訓練集 FAR 分布之涵蓋範圍內（中間 90\% 為 0.54--12.53），故最佳化搜尋未落入分布外推風險。

其二，\textbf{不同路線配置產生可辨識之最佳化差異}：路線 B（開放廣場路線）之最終最佳人行動線熱暴露（$\bar{\Phi}_{\text{walkway}}=0.336$）優於路線 C（$0.341$）與路線 A（$0.342$），且路線 B、C 之最終帕累托前緣規模（均為 $n=40$）亦大於路線 A（$n=20$），顯示當人行路線通過場地中央開放空間時，最佳化引擎有更大的建築配置自由度可同時兼顧陰影遮蔽與路徑熱舒適，而路線 A 之臨街面路線因路徑貼近場地邊界，可調整之建築退縮與遮蔭配置選項相對受限。此一差異證實了本研究之核心論點：人行路線之選址本身即是影響都市設計最佳化結果的關鍵變數，而非與建築量體配置無關之獨立參數。

其三，\textbf{全域 UTCI 目標與路徑選擇之交互作用}：三路線之最佳全域平均 UTCI 差異相對較小（$36.25$--$36.33^\circ\text{C}$），顯示 $f_1$ 對路線選擇之敏感度低於 $f_2$，此一現象合乎預期——$f_1$ 為全場域空間平均，對局部路徑之幾何差異本不敏感，此正是第~\ref{sec:problem}節第二個研究缺口所述「全空間平均最佳化難以反映行人沿路徑之動態熱暴露歷程」之量化實證：若僅以 $f_1$ 作為評估指標，三條路線之設計方案將被判定為近乎等價，唯有透過 $f_2$ 之路徑感知評估，才能區辨出路線 B 之顯著優勢。

\begin{figure}[htbp]
    \centering
    \includegraphics[width=\textwidth]{../urban-thermal-gnn/04_training/figures/fig_abc_comparison.png}
    \caption[A/B/C 三種人行路線配置之三目標 NSGA-II 比較實驗]{A/B/C 三種人行路線配置之三目標 NSGA-II 比較實驗：上排為三條路線於場地上之幾何配置與各自最佳全域 UTCI 解之建築／喬木佈局；下排 (a)(b) 為全域 UTCI 與人行動線熱暴露之收斂曲線，(c) 為三組實驗最終代帕累托前緣之散點對照。}
    \label{fig:abc_comparison}
\end{figure}

需說明之範疇限制：本實驗之路線為固定輸入（非最佳化搜尋對象），驗證的是「已推導之 $\bar{\Phi}_{\text{walkway}}$ 公式與三目標形式化在不同路線幾何下之有效性與區辨力」，而非完整之多路徑 $A^*$ 動線圖同時最佳化；後者仍列為第~\ref{sec:future_work}節之後續工作。


\subsection{設計疊代示範：從帕累托前緣到方案選擇}
\label{subsec:design_iteration_demo}

第~\ref{subsec:abc_comparison}節之三目標最佳化並不輸出單一「最優方案」，而是輸出一組非支配之候選方案集合，實際的設計決策仍須由設計者依當下案址之優先考量從中選擇。圖~\ref{fig:design_iteration} 以路線 B（開放廣場路線）最終帕累托前緣中兩個真實非支配解為例，示範此一比較與選擇過程：方案 1 為前緣上全域平均 UTCI 最低之解（$36.26^\circ\text{C}$），其建築量體集中且 FAR 較高（$1.78$）；方案 2 為前緣上人行動線熱暴露最低之解（$\bar{\Phi}_{\text{walkway}}=0.336$），其建築量體配置退離路線兩側、FAR 較低（$0.71$），以犧牲部分全域平均表現（UTCI 上升至 $36.33^\circ\text{C}$）換取路線熱舒適度之改善。此一比較呈現了本系統作為設計輔助工具之核心互動模式：設計者輸入建築輪廓與人行動線後，工具並非提供單一線性答案，而是提供一組可供權衡比較的候選方案，設計者可依據案址之實際優先權重（例如行人流量較高的路段應優先考慮方案 2 之配置邏輯）進行選擇與後續之量體調整，形成第~\ref{sec:framework}節所述「輸入幾何 → 運算輸出 → 設計回饋調整」之閉環決策過程。

\begin{figure}[htbp]
    \centering
    \includegraphics[width=\textwidth]{../urban-thermal-gnn/04_training/figures/fig_design_iteration.png}
    \caption[設計疊代示範：從建築輪廓輸入到計算輸出、再到設計回饋調整之閉環]{設計疊代示範：以路線 B 帕累托前緣中兩個真實非支配解為例，示範建築量體配置在全域熱舒適與人行動線熱舒適兩種優先考量下之具體權衡差異，模擬設計者於候選方案間比較、選擇並回饋調整之閉環決策情境。}
    \label{fig:design_iteration}
\end{figure}


\section{真實建物擷取流程之案例驗證}
\label{sec:real_case_validation}

第~\ref{subsec:extraction_vs_parametric}節說明了真實建物擷取流程與隨機生成場景之互補分工：前者用於真實都市紋理下之案例驗證，後者支撐模型訓練所需之大規模樣本集。本節以第~\ref{subsec:vectorize_method}節三服務融合方法擷取之真實建物量體，作為 PI-ST-GNN 之推理輸入，檢視模型在未見過真實幾何配置下的表現。

由於本案例擷取之建物幾何（棟距、樓高分布、平面形態）與訓練資料集之 300 個真實場景分布未必完全重疊，此一驗證本質上是對模型一般化能力的分布外（Out-of-Distribution）壓力測試，而非對訓練集內插能力的重複驗證。第~\ref{subsec:coverage}節所述之訓練集幾何特徵覆蓋範圍（FAR 中間 90\% 為 0.54--12.53、BCR 中間 90\% 為 0.18--1.00、簷高 3--162m），可作為判斷真實案例是否落於模型可靠一般化區間的參照基準：若真實場地之量體統計特徵落於此範圍內，可對模型輸出有較高信心；若明顯超出，則預測結果應視為初步參考，建議搭配傳統模擬工具進行交叉驗證。


\section{系統整合架構效能討論}
\label{sec:system_discussion}

\subsection{FastAPI/WebSocket 通訊延遲}
\label{subsec:api_latency}

整合系統的端到端延遲由幾何序列化（Grasshopper → Python）、GNN 推理與結果回傳三個環節組成，合計落在低延遲等級（詳細拆解與各環節毫秒數見附錄~\ref{appx:latency_breakdown}），遠低於人眼察覺延遲（$\sim$1 秒）的門檻，能提供令設計師滿意的「近即時」互動體驗，並支撐第~\ref{subsec:abc_comparison}節所述、涉及數十代族群演化的 NSGA-II 搜尋於可行時程內完成。


\subsection{使用者操作流程概覽}
\label{subsec:user_workflow_overview}

前述效能數字所服務的最終對象，是實際操作 Grasshopper 的設計者；系統是否真正達成即時閉環設計，最終仍須以使用者從「周遭環境模型建好」到「拿到熱舒適評估」之間實際要走的步驟來檢驗，而非僅以延遲毫秒數衡量。圖~\ref{fig:user_workflow} 將本研究實作之 \texttt{UTCIPredictor}（單次即時預測）與 \texttt{UTCIOptimizer}（多目標最佳化搜尋）兩個 Grasshopper 元件之操作流程，統整為單一之 1--8 步驟流程圖：步驟 1--2 為一次性前置準備（啟動背景運算服務、安裝通訊套件）；步驟 3--5 為場地幾何準備與元件連接（匯入基地與建物幾何、接上輸入端口、送出計算請求）；步驟 5 之後依應用情境分岔為路徑 A（單次即時評估，提供立即性回饋、可邊調整設計邊觀察熱力圖）與路徑 B（多目標最佳化搜尋，背景執行緒進行、不阻塞 Rhino 操作、可隨時取消）；步驟 8 為結果讀取與設計回饋，兩路徑之輸出（熱力圖網格、數值指標、帕累托候選方案）皆直接呈現於設計者當下操作之 Rhino 場景中，不另外產生 PDF 報告。

\begin{figure}[htbp]
    \centering
    \includegraphics[width=\textwidth]{../urban-thermal-gnn/04_training/figures/fig_user_workflow.png}
    \caption[Grasshopper 元件操作流程（1--8 步驟）]{Grasshopper 元件操作流程（1--8 步驟）：資料準備（步驟 1--3）→ 匯入元件（步驟 4--5）→ 計算（路徑 A／B，步驟 6--7）→ 結果讀取與設計回饋（步驟 8）。}
    \label{fig:user_workflow}
\end{figure}

完整之逐步操作手冊、輸入輸出端口對照表與常見狀況排除，詳見附錄~\ref{appx:user_manual}。


\subsection{物理一致性與建築設計輔助決策模型的可解釋性}
\label{subsec:trustworthiness}

本系統在建築設計輔助工具的脈絡中，預測結果的\textbf{物理合理性（Physical Plausibility）}比純粹的數值精度更為重要 \citep{shan2025physics}。一個預測「陰影區域比日照區域更熱」的 AI 模型，即便整體 $R^2$ 很高，也可能在關鍵設計決策點上誤導建築師。

物理資訊損失函數在以下具體情境中的功效：

\begin{enumerate}
    \item \textbf{陰影形態評估：}$\mathcal{L}_{\text{rad}}$ 確保模型始終正確反映陰影帶的降溫效果，使建築師在評估不同立面開口率或遮陽板配置時，獲得符合直覺的熱舒適反饋。
    \item \textbf{日間熱積累趨勢：}$\mathcal{L}_{\text{temp}}$ 確保 UTCI 預測曲線符合都市熱慣性的物理規律——早晨緩升、午後達峰、傍晚緩降，使設計師能正確評估不同朝向建築的日間熱積累差異。
    \item \textbf{高層建築風場影響：}$\mathcal{L}_{\text{wind}}$ 確保高層建築背風側的「熱死角」能被正確識別，這是建築設計中常被忽略但對行人熱舒適影響顯著的空氣動力學現象 \citep{oke1987boundary}。
\end{enumerate}

上述物理一致性的維持，使得 PI-ST-GNN 不僅是一個「黑箱預測器」，更是一個\textbf{具有熱力學解釋能力的設計輔助工具}，能在設計師的專業判斷基礎上提供可驗證的科學依據。

為進一步驗證此一物理一致性主張並非僅為訓練指標上的巧合，本研究以 GNNExplainer 對本研究最終採用之 V5-300 模型進行事後可解釋性分析 \citep{shan2025physics}，結果如圖~\ref{fig:gnnexplainer} 所示。以真實場景資料集中一處熱壓力峰值節點（$\text{UTCI}=41.7°\text{C}$）為分析目標，梯度顯著性分析顯示該節點之預測值主要由其空間鄰域中的空氣溫度 $T_a$（正規化歸因值 1.00）主導，平均輻射溫度 $T_{\text{mrt}}$ 居次（0.79），相對濕度 RH 居第三（0.53），陰影旗標 $f_{sh}$、風速 $v_a$ 與天空視角因子 SVF 則為次要貢獻，此一特徵重要性排序反映真實場景中氣溫與輻射環境對熱壓力之主導地位，與都市微氣候學之物理直覺一致。

RGCN 各層之關係權重范數分析（附錄~\ref{appx:dynamic_edges_fix}記錄之真實場景資料集動態邊修正，五類關係於訓練時皆具真實動態邊）顯示，五類關係之權重范數皆落於相近量級（見圖~\ref{fig:gnnexplainer}(c)），未呈現單一關係獨大之型態，且\textbf{不再出現任何「無真實訓練邊」之虛線填充樣式}——此為 V5-300 動態邊缺陷修正後之直接可視化證據：淺層（第 1--2 層）中各關係權重相近，反映模型在淺層同時整合多種物理管道之訊號；深層則呈現 \texttt{contiguity}（空間毗鄰）權重上升之分化趨勢，與空間毗鄰關係需經多跳擴散才能涵蓋建築陰影最大投射距離（第~\ref{subsec:rgcn}節）之設計動機相符。此一分佈型態較單一關係獨大更符合物理直覺：陰影、植被蒸散、對流、語義、毗鄰五種因果機制理應共同、而非由單一路徑主導地貢獻於熱壓力預測。此一可解釋性驗證，補強了單純數值指標難以呈現的「模型是否學到正確物理機制」此一問題的實證依據。

\begin{figure}[htbp]
    \centering
    \includegraphics[width=\textwidth]{../urban-thermal-gnn/04_training/figures/fig_gnnexplainer.png}
    \caption[GNNExplainer 可解釋性分析]{PI-ST-GNN 之 GNNExplainer 事後可解釋性分析（真實場景資料集場景 4，V5-300 checkpoint，驗證 $R^2=0.9837$）}
    \label{fig:gnnexplainer}
\end{figure}

上圖 (a) 為目標節點（黑圈標示，UTCI=41.7°C）周邊之空間梯度顯著性，顏色越深代表該節點對目標預測值之影響越大；(b) 為九維輸入特徵之歸因排序，空氣溫度為主導特徵；(c) 為三層 RGCN 各關係類型之權重 Frobenius 范數，五類關係皆具真實訓練邊、范數量級相近；(d) 為目標節點兩跳鄰域內前 40 條重要邊之子圖，邊寬與顏色深淺代表梯度歸因強度。


\subsection{系統局限性與誤差來源}
\label{subsec:limitations}

儘管本系統取得優異的整體預測性能，仍存在若干局限性值得說明：

\begin{enumerate}
    \item \textbf{氣候情境涵蓋範圍之限制：}本研究另一條程序化主線之訓練資料集僅涵蓋夏至（6月21日）之單一典型氣候日；第~\ref{sec:v4_validation}節之主要報告模型（V5-300）真實資料集雖已改以 2025 年 6、7、8 月之真實逐時觀測氣象驅動，將氣候情境由單日典型值擴展為真實夏季逐時變異，但仍未涵蓋冬季東北季風、梅雨季陰天等其他季節條件，模型對此類條件之一般化能力尚待驗證；若應用於其他季節，建議以對應季節之真實觀測或 EPW 重新微調。
    \item \textbf{平流與濕球效應未建模：}現有模型未明確建模都市熱島效應中的大尺度平流傳熱，以及濕地、噴泉等蒸發冷卻設施的局部降溫效果，此類效應在極端熱壓力情境下可能不可忽略。
    \item \textbf{極端熱壓力等級準確率偏低：}如消融實驗所示，極端熱壓力（UTCI $> 46°\text{C}$）之分類準確率明顯低於其他等級，主因是訓練資料中此等級樣本佔比極低（僅出現於正午直射、高 SVF 條件下），建議後續研究採用加權採樣或針對高溫情境進行額外資料擴增。
    \item \textbf{單一行區尺度限制：}本研究聚焦於單一 $80 \times 80$ 公尺行區，尚未考量相鄰行區間的微氣候互動（如建築群之間形成的峽谷效應或大尺度通風廊道），多行區聯合建模為未來重要研究方向。
    \item \textbf{真實建物擷取之樓高辨識限制：}第~\ref{subsec:vectorize_method}節之光學文字辨識於複合式大型建物之並列樓高分區仍偶有辨識失敗或分區未完全切割之情形（已於流程中如實標記為待覆核案例，而非以猜測值填補），此類案例於本章之真實案例驗證中須人工複核樓高屬性後方可採用。
\end{enumerate}
